% Dunnhumby data understanding and preparation, adapted from:
% https://github.com/deepmancer/latex-presentation-template (MIT, 2024 Alireza Heidari).
% Compile in this directory: xelatex -interaction=nonstopmode -halt-on-error main.tex
% Compile twice to update the slide counter. Set Overleaf compiler to XeLaTeX.
\documentclass[smaller,aspectratio=169,t]{beamer}
\usepackage{fontspec}
\IfFontExistsTF{Noto Sans}{\setsansfont{Noto Sans}}{%
  \IfFontExistsTF{Arial}{\setsansfont{Arial}}{\setsansfont{TeX Gyre Heros}}}
\usepackage{tabularx,array,graphicx,tikz}
\IfFileExists{booktabs.sty}{\usepackage{booktabs}}{%
  \newcommand{\toprule}{\hline}
  \newcommand{\midrule}{\hline}
  \newcommand{\bottomrule}{\hline}}
\usetikzlibrary{arrows.meta,positioning}
\definecolor{myBlue}{RGB}{0,33,71}
\definecolor{accent}{HTML}{008B8B}
\definecolor{pale}{HTML}{EEF4F7}
\definecolor{muted}{HTML}{5D6D7E}
\useoutertheme{infolines}
\usecolortheme[named=myBlue]{structure}
\setbeamertemplate{navigation symbols}{}
\setbeamertemplate{headline}{}
\setbeamersize{text margin left=10mm,text margin right=10mm}
\setbeamercolor{normal text}{fg=myBlue,bg=white}
\setbeamercolor{frametitle}{fg=myBlue,bg=white}
\setbeamerfont{frametitle}{size=\large,series=\bfseries}
\setbeamercolor{block title}{fg=white,bg=myBlue}
\setbeamercolor{block body}{fg=myBlue,bg=pale}
\setbeamertemplate{frametitle}{%
  \nointerlineskip
  \begin{beamercolorbox}[wd=\paperwidth,leftskip=10mm,rightskip=10mm,
    ht=8mm,dp=2mm]{frametitle}%
    \usebeamerfont{frametitle}\insertframetitle
  \end{beamercolorbox}%
  \nointerlineskip\hbox to \paperwidth{\hspace{10mm}%
    \color{myBlue}\rule{\dimexpr\paperwidth-20mm\relax}{.35pt}\hfil}%
  \vspace{1mm}}
\setbeamertemplate{footline}{%
  \leavevmode\hbox{%
    \begin{beamercolorbox}[wd=.64\paperwidth,ht=3.2mm,dp=1.6mm,leftskip=10mm]{footline}%
      \tiny\color{muted}Dunnhumby | Data Understanding \& EDA
    \end{beamercolorbox}%
    \begin{beamercolorbox}[wd=.25\paperwidth,ht=3.2mm,dp=1.6mm,center]{footline}%
      \tiny\color{muted}[Tên người trình bày]
    \end{beamercolorbox}%
    \begin{beamercolorbox}[wd=.11\paperwidth,ht=3.2mm,dp=1.6mm,center]{footline}%
      \tiny\color{muted}\insertframenumber/\inserttotalframenumber
    \end{beamercolorbox}}\vspace{1mm}}
\setbeamertemplate{itemize item}{\color{accent}$\blacktriangleright$}
\newcommand{\src}[1]{\vspace{1mm}{\tiny\color{muted}Nguồn: #1}}
\newcommand{\fig}[2]{\includegraphics[width=#1,height=.41\textheight,keepaspectratio]{figures/#2.pdf}}
\newcommand{\largefig}[1]{\includegraphics[width=.96\linewidth,height=.58\textheight,keepaspectratio]{figures/#1.pdf}}
\newcommand{\edafig}[1]{\includegraphics[width=.98\linewidth,height=.57\textheight,keepaspectratio]{eda_figures/#1.png}}
\newcommand{\edashort}[1]{\includegraphics[width=.98\linewidth,height=.53\textheight,keepaspectratio]{eda_figures/#1.png}}
\newcommand{\kpi}[2]{\begin{minipage}{.235\textwidth}\centering
  {\LARGE\bfseries\color{myBlue}#1}\\[2mm]{\small #2}\end{minipage}}
\title{Phân tích dữ liệu Dunnhumby cho bài toán dự đoán mua hàng}
\subtitle{Dữ liệu có gì, đáng tin đến đâu và hành vi mua nói lên điều gì?}
\author{[Tên người trình bày]}
\date{[Ngày báo cáo]}

\begin{document}

% 01
\begin{frame}[plain]
\begin{tikzpicture}[remember picture,overlay]
  \fill[myBlue] (current page.south west) rectangle (current page.north east);
  \fill[accent] ([yshift=14mm]current page.south west) rectangle
    ([yshift=16mm]current page.south east);
  \fill[white,opacity=.05] ([xshift=-50mm]current page.north east) --
    (current page.north east) -- ([yshift=35mm]current page.south east) -- cycle;
\end{tikzpicture}
\vspace{6mm}
{\color{white}\LARGE\bfseries PHÂN TÍCH DỮ LIỆU DUNNHUMBY\\[2mm]
CHO BÀI TOÁN DỰ ĐOÁN MUA HÀNG}\par
\vspace{6mm}{\color{white}\large Giai đoạn tìm hiểu và khám phá dữ liệu đã hoàn thành}\par
\vspace{12mm}{\color{white}[Tên người trình bày] \quad | \quad [Ngày báo cáo]}
\end{frame}

% 02
\begin{frame}{Mục tiêu và phạm vi nghiên cứu}
\begin{columns}[T,onlytextwidth]
\column{.49\textwidth}
\begin{block}{Câu hỏi cần trả lời}
\begin{itemize}
\item Bộ dữ liệu có những bảng nào, liên kết với nhau ra sao?
\item Một dòng giao dịch, một giỏ hàng và một hộ khác nhau thế nào?
\item Dữ liệu có thiếu, trùng hoặc giá trị khó giải thích không?
\item Hành vi trước đây liên quan thế nào đến việc mua trong 7 ngày tới?
\end{itemize}
\end{block}
\column{.47\textwidth}
\textbf{Lộ trình}\vspace{3mm}
\begin{enumerate}
\item Dữ liệu gồm những gì?
\item Chất lượng và cách kết nối các bảng
\item Hành vi mua của hộ gia đình
\item Những điều đã biết và còn chưa chắc chắn
\end{enumerate}
\vspace{4mm}
{\small Phần này chỉ trình bày dữ liệu và kết quả khám phá, chưa đánh giá mô hình.}
\end{columns}
\end{frame}

% 03
\begin{frame}{Tổng quan quy mô bộ dữ liệu}
\vspace{4mm}
\kpi{2.500}{hộ gia đình}
\kpi{276.484}{giỏ hàng}
\kpi{2.595.732}{dòng sản phẩm}
\kpi{711}{ngày dữ liệu}
\vspace{9mm}
\begin{block}{Bối cảnh}
Đây là lịch sử mua sắm của 2.500 hộ tại một nhà bán lẻ, trải dài khoảng hai năm.
Ngoài giao dịch, dữ liệu còn có sản phẩm, thông tin hộ, chiến dịch, phiếu giảm giá và trưng bày.
\end{block}
\src{Audit local; \href{https://www.dunnhumby.com/source-files/}{dunnhumby Source Files --- The Complete Journey}.}
\end{frame}

% 04
\begin{frame}{Danh mục tám bảng dữ liệu}
\small\renewcommand{\arraystretch}{1.08}
\begin{tabularx}{\textwidth}{>{\ttfamily}l r X}
\toprule Bảng & Số dòng & Mỗi dòng thể hiện \\
\midrule
transaction\_data & 2.595.732 & Một sản phẩm trong một giỏ hàng \\
product & 92.353 & Một sản phẩm và nhóm hàng của nó \\
hh\_demographic & 801 & Thông tin mô tả một hộ gia đình \\
campaign\_table & 7.208 & Một hộ được gắn với một chiến dịch \\
campaign\_desc & 30 & Một chiến dịch và thời gian diễn ra \\
coupon & 124.548 & Một liên kết phiếu giảm giá -- sản phẩm -- chiến dịch \\
coupon\_redempt & 2.318 & Một lần sử dụng phiếu giảm giá \\
causal\_data & 36.786.524 & Một ghi nhận trưng bày/gửi quảng cáo theo sản phẩm -- cửa hàng -- tuần \\
\bottomrule
\end{tabularx}
\src{Notebook phân tích dữ liệu, mục 1; tên bảng causal\_data không có nghĩa dữ liệu chứng minh quan hệ nhân quả.}
\end{frame}

% 05
\begin{frame}{Một hộ có thể có nhiều giỏ hàng và nhiều dòng sản phẩm}
\begin{center}
\begin{tikzpicture}[node distance=9mm,
  box/.style={draw=myBlue,rounded corners,minimum height=10mm,text width=33mm,align=center,fill=pale},
  arrow/.style={-{Latex},thick,color=accent}]
\node[box] (h) {1 hộ gia đình};
\node[box,right=of h] (b) {N giỏ hàng};
\node[box,right=of b] (l) {N dòng sản phẩm};
\draw[arrow] (h)--(b); \draw[arrow] (b)--(l);
\end{tikzpicture}
\end{center}
\vspace{8mm}
\begin{columns}[T,onlytextwidth]
\column{.48\textwidth}
\textbf{Ví dụ cấu trúc}\vspace{1mm}
\begin{tabular}{rrr}\toprule hộ & giỏ & sản phẩm \\
1 & A & sữa \\ 1 & A & bánh \\ 1 & B & gạo \\
\bottomrule\end{tabular}
\column{.48\textwidth}
\begin{itemize}
\item 3 dòng sản phẩm thuộc 2 giỏ của 1 hộ.
\item Đếm giỏ theo số mã \texttt{BASKET\_ID} khác nhau.
\item Một hộ có thể có nhiều giỏ trong cùng ngày.
\end{itemize}
\end{columns}
\end{frame}

% 06
\begin{frame}{Cấu trúc liên kết giữa các bảng dữ liệu}
\centering
\begin{tikzpicture}[scale=.78,transform shape,node distance=8mm,
  box/.style={draw=myBlue,rounded corners,minimum height=9mm,text width=30mm,align=center,fill=pale,font=\small},
  arrow/.style={-{Latex},thick,color=accent}]
\node[box,fill=myBlue,text=white] (tx) {Giao dịch};
\node[box,left=15mm of tx] (hh) {Hộ gia đình};
\node[box,right=15mm of tx] (pr) {Sản phẩm};
\node[box,above=9mm of hh] (demo) {Thông tin hộ};
\node[box,below=9mm of hh] (camp) {Chiến dịch};
\node[box,below=9mm of pr] (coup) {Phiếu giảm giá};
\node[box,above=9mm of pr] (med) {Trưng bày / mailer};
\draw[arrow] (hh)--node[above,font=\tiny]{household\_key}(tx);
\draw[arrow] (pr)--node[above,font=\tiny]{PRODUCT\_ID}(tx);
\draw[arrow] (demo)--(hh); \draw[arrow] (camp)--(hh);
\draw[arrow] (coup)--(pr); \draw[arrow] (med)--(pr);
\end{tikzpicture}
\vspace{3mm}
\begin{block}{Hàm ý đối với kết nối dữ liệu}
Một phiếu giảm giá có thể gắn với nhiều sản phẩm; thông tin trưng bày được ghi theo sản phẩm, cửa hàng và tuần.
Khi nối bảng, cần kiểm tra số dòng để tránh vô tình nhân đôi giao dịch.
\end{block}
\end{frame}

% 07
\begin{frame}{Dữ liệu được kiểm tra và chuẩn bị ra sao?}
\begin{enumerate}
\item \textbf{Kiểm kê:} đếm dòng, cột và xác định mỗi dòng đại diện cho điều gì.
\item \textbf{Chuẩn hóa cách đọc:} nhận diện mã, ngày, tiền; không sửa trực tiếp dữ liệu gốc.
\item \textbf{Kiểm tra chất lượng:} tìm dữ liệu thiếu, dòng trùng và giá trị bất thường.
\item \textbf{Nối bảng:} kiểm tra khóa và số dòng trước, sau khi nối.
\item \textbf{Tổng hợp:} từ dòng sản phẩm $\rightarrow$ giỏ hàng $\rightarrow$ hoạt động của hộ theo ngày.
\item \textbf{Khám phá:} xem nhịp mua, nhóm hàng, mức độ bao phủ và so sánh các nhóm hộ.
\end{enumerate}
\vspace{3mm}\begin{block}{Kết quả đầu ra}Mỗi bảng thống kê đều nêu rõ đang đếm dòng, giỏ hàng hay hộ gia đình.\end{block}
\end{frame}

% 08
\begin{frame}{Kiểm tra dữ liệu thiếu và độ bao phủ}
\begin{center}\largefig{household_coverage}\end{center}
\vspace{-2mm}
\begin{itemize}\small
\item \textbf{Kết quả:} bảng giao dịch không có ô trống; chỉ 801/2.500 hộ có thông tin nhân khẩu học.
\item \textbf{Cách xử lý:} khi nối bảng thông tin hộ, vẫn giữ cả những hộ không có thông tin này.
\end{itemize}
\src{Kiểm tra dữ liệu gốc; mẫu số là 2.500 hộ có giao dịch.}
\end{frame}

% 09
\begin{frame}{Kiểm tra dữ liệu trùng lặp}
\begin{columns}[T,onlytextwidth]
\column{.50\textwidth}
\begin{block}{Phát hiện}
\begin{itemize}
\item Giao dịch: không có hai dòng giống hệt nhau trên cả 12 cột.
\item Sản phẩm: mỗi \texttt{PRODUCT\_ID} chỉ xuất hiện một lần trong danh mục.
\item Phiếu giảm giá: có 5.164 dòng giống hệt nhau trên cả 3 cột.
\end{itemize}
\end{block}
\column{.47\textwidth}
\textbf{Hàm ý xử lý}\vspace{2mm}
\begin{itemize}
\item Giữ bản gốc và ghi nhận số dòng trùng.
\item Chỉ loại dòng phiếu giảm giá trùng sau khi xác định đúng ý nghĩa mỗi dòng.
\item Không coi nhiều sản phẩm trong cùng giỏ là dòng trùng.
\end{itemize}
\end{columns}
\end{frame}

% 10
\begin{frame}{Kiểm tra tính toàn vẹn trước khi kết nối dữ liệu}
\small\renewcommand{\arraystretch}{1.2}
\begin{tabularx}{\textwidth}{Xrl}
\toprule Nội dung kiểm tra & Kết quả & Ý nghĩa \\
\midrule
Sản phẩm giao dịch có trong danh mục & 100\% & Có thể nối theo mã sản phẩm \\
Một mã giỏ thuộc nhiều hộ & 0 & Mỗi giỏ thuộc đúng một hộ \\
Một mã giỏ thuộc nhiều ngày & 0 & Ngày ghi nhận nhất quán \\
Lần dùng phiếu có mã chiến dịch hợp lệ & 100\% & Mã có trong bảng chiến dịch \\
Lần dùng phiếu nằm trong thời gian chiến dịch & 100\% & Ngày dùng nằm trong khoảng được ghi \\
\bottomrule
\end{tabularx}
\vspace{4mm}
\begin{block}{Hàm ý phân tích}
Mã sản phẩm không trùng trong danh mục và mọi giao dịch đều tìm được sản phẩm tương ứng.
Vì vậy có thể nối bảng sản phẩm mà không làm tăng số dòng giao dịch.
\end{block}
\src{Kiểm tra trên các bảng được cung cấp; không xác nhận nội dung thực tế của chiến dịch.}
\end{frame}

% 11
\begin{frame}{QUANTITY: biết lượng mua, chưa rõ đơn vị ở mọi nhóm hàng}
\begin{center}\largefig{quantity_semantics}\end{center}
\vspace{-2mm}
\begin{itemize}\small
\item Có 23.136 dòng \texttt{QUANTITY > 100}, chủ yếu ở nhiên liệu và nhóm hàng khác.
\item Tài liệu chỉ gọi đây là lượng mua, không nêu rõ đơn vị riêng cho nhiên liệu; vì vậy chưa cộng lượng mua giữa các nhóm hàng.
\end{itemize}
\src{Audit transaction--product; The Complete Journey User Guide, trang 3.}
\end{frame}

% 12
\begin{frame}{Giỏ hàng có doanh số bằng 0 vẫn là lượt mua được ghi nhận}
\begin{columns}[T,onlytextwidth]
\column{.48\textwidth}
{\LARGE\bfseries 942}\\[1mm]giỏ có tổng doanh số bằng 0\vspace{4mm}
\begin{itemize}
\item 275.542 giỏ có tổng doanh số dương.
\item 18.850 dòng sản phẩm có doanh số bằng 0.
\item Dòng 0 đồng chưa chắc là lỗi; dữ liệu không cho biết đầy đủ nguyên nhân.
\end{itemize}
\column{.48\textwidth}
\begin{block}{Định nghĩa đã chốt}
Nhãn cần dự đoán là \textbf{có bất kỳ giỏ hàng nào} trong 7 ngày tới; kể cả giỏ 0 đồng, nhiên liệu và nhóm hàng khác.
\end{block}
\vspace{2mm}
\textbf{Vì vậy:} không loại giỏ hàng chỉ vì doanh số bằng 0.
\end{columns}
\end{frame}

% 13
\begin{frame}{Phân bố số giỏ hàng theo ngày}
\begin{center}\largefig{daily_baskets}\end{center}
\vspace{-2mm}\begin{itemize}\small
\item Số giỏ tăng ở giai đoạn đầu; riêng ngày 278 và 643 chỉ ghi nhận 1 giỏ/ngày.
\item Chưa biết vì sao hai ngày này thấp bất thường; đường đậm thể hiện trung bình 14 ngày.
\end{itemize}
\src{Bảng giao dịch; \texttt{DAY} là số thứ tự ngày, chưa xác định được ngày lịch.}
\end{frame}

% 14
\begin{frame}{Phân bố thời điểm xuất hiện đầu tiên của hộ}
\begin{center}\largefig{first_observed}\end{center}
\vspace{-2mm}\begin{itemize}\small
\item Một nửa số hộ bắt đầu xuất hiện từ ngày \textbf{69} trở về trước.
\item Đây là ngày mua đầu tiên \textbf{được ghi trong dữ liệu}, không chắc là ngày họ bắt đầu làm khách hàng.
\end{itemize}
\end{frame}

% 15
\begin{frame}{Phân bố số giỏ hàng theo hộ gia đình}
\begin{center}\largefig{basket_per_household}\end{center}
\vspace{-2mm}\begin{itemize}\small
\item Một hộ có trung vị \textbf{79} giỏ, trung bình \textbf{110,6} giỏ trong toàn kỳ.
\item Một số hộ mua nhiều hơn hẳn phần còn lại; có 39.179 lượt hộ--ngày với từ 2 giỏ trở lên.
\end{itemize}
\src{Biểu đồ chỉ hiển thị đến 600 giỏ/hộ; số liệu thống kê vẫn tính trên toàn bộ dữ liệu.}
\end{frame}

% 16
\begin{frame}{Phân bố số dòng sản phẩm trong giỏ hàng}
\begin{center}\largefig{lines_per_basket}\end{center}
\vspace{-2mm}\begin{itemize}\small
\item Một giỏ có trung vị \textbf{5 dòng sản phẩm}, trung bình \textbf{9,39 dòng}.
\item Có những giỏ lớn hơn nhiều so với giỏ thông thường; không nhầm số dòng với số giỏ.
\end{itemize}
\src{Biểu đồ chỉ hiển thị đến 60 dòng/giỏ; số liệu thống kê vẫn tính trên toàn bộ dữ liệu.}
\end{frame}

% 17
\begin{frame}{Cơ cấu doanh số ghi nhận theo nhóm hàng}
\begin{center}\largefig{department_sales}\end{center}
\vspace{-2mm}\begin{itemize}\small
\item Nhóm tạp hóa (GROCERY) chiếm \textbf{50,8\%} doanh số; DRUG GM 13,1\%; rau quả (PRODUCE) 6,9\%.
\item Dữ liệu không chỉ có tạp hóa mà còn ghi nhận nhiên liệu và các nhóm hàng khác.
\end{itemize}
\src{Tổng SALES\_VALUE retailer nhận được; hiển thị 8 department lớn nhất. User Guide, trang 3.}
\end{frame}

% 18
\begin{frame}{Cơ cấu doanh số ghi nhận theo loại nhãn hàng}
\begin{center}\largefig{brand_mix}\end{center}
\vspace{-2mm}\begin{itemize}\small
\item Nhãn hàng của nhà sản xuất (National) chiếm \textbf{72,2\%} doanh số; nhãn riêng (Private) chiếm \textbf{27,8\%}.
\item Tỷ lệ này tính theo doanh số của các hộ trong bộ dữ liệu, không đại diện toàn bộ khách hàng của nhà bán lẻ.
\end{itemize}
\end{frame}

% 19
\begin{frame}{Dữ liệu bổ sung về chiến dịch và thông tin hộ}
\small
\begin{tabularx}{\textwidth}{>{\bfseries}p{3cm}X X}
\toprule Nguồn & Quy mô/độ phủ & Mỗi dòng thể hiện \\
\midrule
Chiến dịch & 63,36\% hộ được gắn chiến dịch & Một hộ -- một chiến dịch. \\
Sử dụng phiếu & 2.318 lần sử dụng & Một hộ -- một phiếu -- một ngày. \\
Trưng bày / mailer & 36,79 triệu dòng & Một sản phẩm -- cửa hàng -- tuần. \\
Thông tin hộ & 801 hộ có hồ sơ & Một hộ; không bao phủ cả mẫu. \\
\bottomrule
\end{tabularx}
\vspace{4mm}\begin{block}{Hàm ý nghiên cứu}
Những bảng này giúp hiểu thêm bối cảnh mua sắm, nhưng không cho biết chắc khách đã nhìn thấy quảng cáo hay vì sao họ mua hàng.
\end{block}
\end{frame}

% EDA 01
\begin{frame}{Thế nào được tính là mua hàng trong 7 ngày tới?}
\begin{center}\edafig{target_scenarios}\end{center}
\vspace{-2mm}\begin{itemize}\small
\item Một hộ được tính là \textbf{có mua} nếu xuất hiện ít nhất một mã giỏ trong 7 ngày sau mốc chốt $s$. Tỷ lệ này là 52,79\% trên 129.834 lần quan sát hộ--mốc.
\item Nếu chỉ tính giỏ có doanh số dương, 33 nhãn đổi; nếu chỉ tính tạp hóa, 5.309 nhãn đổi.
\end{itemize}
\src{Notebook phân tích dữ liệu, mục 5.2; hai cách tính thay thế chỉ dùng để kiểm tra định nghĩa.}
\end{frame}

% EDA 02
\begin{frame}{Vì sao cần mốc chốt dữ liệu cho từng hộ?}
\vspace{3mm}
\begin{center}
\begin{tikzpicture}[x=1cm,y=1cm,>=Latex]
\draw[very thick,accent,-{Latex}] (0,0)--(10.4,0);
\draw[fill=pale,draw=myBlue] (.4,-.35) rectangle (6.0,.35);
\draw[fill=accent!20,draw=accent] (6.0,-.35) rectangle (9.2,.35);
\draw[thick] (6.0,-.5)--(6.0,.55);
\node[above] at (3.2,.35) {Lịch sử quan sát đến $s$};
\node[above] at (7.6,.35) {Nhãn $(s,s+7]$};
\node[below] at (6.0,-.5) {mốc chốt $s$};
\end{tikzpicture}
\end{center}
\vspace{5mm}
\begin{itemize}
\item Mỗi lần quan sát gồm \textbf{một hộ tại một mốc chốt}; cùng một hộ có thể được quan sát ở nhiều tuần.
\item Chỉ dùng thông tin đã có đến mốc chốt để xem hộ có mua trong 7 ngày sau đó không.
\item Các mốc dùng để phân tích trải từ ngày 180 đến 537, cách nhau 7 ngày; hộ cần từng có giỏ hàng trước hoặc tại mốc.
\item Vì một hộ xuất hiện nhiều lần, tỷ lệ theo số lần quan sát không phải tỷ lệ theo số hộ khác nhau.
\end{itemize}
\src{Notebook EDA, mục 5; \texttt{DAY} là chỉ số ngày tương đối.}
\end{frame}

% EDA 03
\begin{frame}{Mua gần đây liên quan đến khả năng mua tiếp}
\begin{center}\edafig{recency_activity}\end{center}
\vspace{-2mm}\begin{itemize}\small
\item Trong nhóm vừa mua cách đây 0--7 ngày, 70,47\% có giỏ trong 7 ngày tiếp theo; nhóm đã hơn 90 ngày chưa mua chỉ đạt 6,56\%.
\item Đếm \textbf{số ngày có mua} trong 90 ngày để đo nhịp hoạt động; một ngày có thể có nhiều giỏ.
\end{itemize}
\src{Notebook phân tích dữ liệu, mục 6; đây là mối liên hệ quan sát được, không chứng minh nguyên nhân.}
\end{frame}

% EDA 04
\begin{frame}{Chọn 40 hộ điển hình để nhìn rõ nhịp mua}
\vspace{3mm}
\begin{block}{Cách chọn hộ minh họa tại mốc ngày 537}
Chia hộ theo số \textbf{ngày có mua trong 90 ngày}: ít nhất, gần trung bình, gần trung vị và nhiều nhất. Mỗi nhóm chọn 10 hộ khác nhau.
\end{block}
\vspace{4mm}
\begin{itemize}
\item Trung bình và trung vị được tính trên toàn bộ hộ đủ điều kiện tại mốc chốt, theo số ngày mua chứ không theo số giỏ.
\item Ô màu trong lịch cho biết số giỏ mỗi ngày; trục tuần chỉ là tuần tương đối, không xác định thứ trong tuần.
\item Các hộ được chọn để \textbf{minh họa bốn mức hoạt động}, không phải mẫu đại diện cho toàn bộ dữ liệu.
\end{itemize}
\vspace{3mm}
\src{Notebook EDA, mục 6.1, chứa đủ bốn biểu đồ lịch hoạt động.}
\end{frame}

% EDA 05
\begin{frame}{Mua ở nhiều nhóm hàng đi kèm tỷ lệ mua tiếp cao hơn}
\begin{center}\edafig{department_activity}\end{center}
\vspace{-2mm}\begin{itemize}\small
\item Xét toàn bộ mẫu, hộ mua ở 1--3 nhóm hàng có tỷ lệ mua tiếp 23,13\%; hộ mua ở ít nhất 11 nhóm đạt 84,75\%.
\item Nếu cùng có 4--9 ngày mua trong 90 ngày, tỷ lệ tương ứng là 33,31\% và 51,56\%.
\end{itemize}
\src{Notebook phân tích dữ liệu, mục 6.2; nhóm dưới 30 lượt quan sát hộ được tô xám.}
\end{frame}

% EDA 06
\begin{frame}{Giỏ có giảm giá không đồng nghĩa mua tiếp nhiều hơn}
\begin{center}\edashort{discount_activity}\end{center}
\vspace{-2mm}\begin{itemize}\small
\item Một giỏ được tính là có giảm giá nếu ít nhất một dòng sản phẩm ghi nhận giá trị giảm ở một trong ba cột liên quan.
\item Ở 90.618 lần quan sát hộ, hơn 75\% giỏ trong 90 ngày từng có giảm giá. Đây là \textbf{tỷ lệ giỏ}, không phải tiết kiệm 75\% tiền.
\end{itemize}
\src{Notebook phân tích dữ liệu, mục 6.3; chỉ biết giảm giá tại giao dịch, không biết hộ đã nhìn thấy ưu đãi nào.}
\end{frame}

% EDA 07
\begin{frame}{Nhiên liệu: chênh lệch giảm mạnh khi so hộ có nhịp mua giống nhau}
\begin{center}\edafig{fuel_activity}\end{center}
\vspace{-2mm}\begin{itemize}\small
\item Nếu không chia nhóm, hộ từng mua nhiên liệu có tỷ lệ mua tiếp 73,39\%, so với 45,05\% ở hộ không mua.
\item Nhưng trong nhóm cùng có 4--9 ngày mua, hai tỷ lệ chỉ còn 41,89\% và 40,93\%.
\end{itemize}
\src{Notebook phân tích dữ liệu, mục 7; so cùng nhịp mua vẫn chưa kiểm soát được mọi khác biệt giữa các hộ.}
\end{frame}

% EDA 08
\begin{frame}{Phiếu giảm giá và thông tin hộ: cần xem số mẫu trước khi kết luận}
\vspace{2mm}\small\renewcommand{\arraystretch}{1.28}
\begin{tabularx}{\textwidth}{p{29mm}X X}
\toprule Chủ đề & Điều quan sát được & Cần lưu ý \\
\midrule
Sử dụng phiếu & Nhóm dùng phiếu và chỉ mua 1--3 ngày: 9 lần quan sát, đều từ 1 hộ. & Nhóm quá nhỏ để đại diện; dùng phiếu khác với việc chỉ nhìn thấy ưu đãi. \\
Thông tin hộ & 801/2.500 hộ có hồ sơ; tỷ lệ mua tiếp của nhóm có/không có hồ sơ là 76,61\%/41,55\%. & Nhóm có hồ sơ có thể khác về nhịp mua; chưa rõ hồ sơ có từ thời điểm nào. \\
\bottomrule
\end{tabularx}
\vspace{5mm}
\begin{block}{Quy tắc diễn giải}
Không thể kết luận phiếu giảm giá làm khách mua thêm. Cũng chưa nên dùng thông tin hộ cho các mốc quá khứ khi chưa biết lúc đó thông tin đã tồn tại hay chưa.
\end{block}
\src{Notebook phân tích dữ liệu, mục 8--9; cùng một hộ có thể được tính ở nhiều mốc.}
\end{frame}

% EDA 09
\begin{frame}{Theo thời gian, cần đọc tỷ lệ mua cùng số hộ trong nhóm}
\begin{center}\edafig{activity_over_time}\end{center}
\vspace{-2mm}\begin{itemize}\small
\item Đường tỷ lệ của nhóm ít hộ dễ dao động, nên cần nhìn cả biểu đồ số lượt quan sát bên dưới.
\item Mốc chốt 271 được đánh dấu vì 7 ngày theo sau có ngày 278 chỉ ghi nhận 1 giỏ.
\end{itemize}
\src{Notebook EDA, mục 6.5; biểu đồ không xác định nguyên nhân biến động.}
\end{frame}

% EDA 10
\begin{frame}{Hai ngày có rất ít giỏ: ảnh hưởng đến phần nào của phân tích?}
\vspace{3mm}\small\renewcommand{\arraystretch}{1.3}
\begin{tabularx}{\textwidth}{lXr}
\toprule Ngày dữ liệu & Mốc chốt có 7 ngày tiếp theo chứa ngày này & Số hộ đủ điều kiện \\
\midrule
278 & Mốc 271, phần dùng để phân tích & 2.496 \\
643 & Mốc 642, phần để đánh giá sau này & 2.499 \\
\bottomrule
\end{tabularx}
\vspace{5mm}
\begin{itemize}
\item Cửa sổ 90 ngày trước mốc của 13 lần quan sát ở phần phân tích và 8 lần ở phần đánh giá cũng chứa một trong hai ngày.
\item Đây là \textbf{phạm vi có thể bị ảnh hưởng}, không phải số nhãn bị sai.
\item Chưa biết ngày lịch hay nguyên nhân; giữ dữ liệu gốc và đánh dấu để kiểm tra thêm.
\end{itemize}
\src{Notebook EDA, mục 10; phép đếm dựa trên thời gian và eligibility, không dùng nhãn test.}
\end{frame}

% EDA 11
\begin{frame}{RFM: hộ mua gần đây, thường xuyên và bao nhiêu tiền?}
\begin{center}\edashort{rfm_distributions}\end{center}
\vspace{-2mm}\begin{itemize}\small
\item Tại mốc ngày 537: $R$ là số ngày từ lần mua gần nhất; $F$ là số giỏ trong 90 ngày; $M$ là tổng doanh số trong 90 ngày.
\item Mỗi hộ xuất hiện một lần. Trên biểu đồ, giá trị vượt phân vị 99 được gom vào cột cuối để dễ nhìn; không xóa chúng khỏi dữ liệu.
\end{itemize}
\src{Notebook phân tích dữ liệu, mục 11.1; $M$ là doanh số ghi nhận, không phải lợi nhuận hay giá trị vòng đời khách hàng.}
\end{frame}

% EDA 12
\begin{frame}{Mua gần và mua thường xuyên: đọc hai chiều cùng nhau}
\begin{center}\edafig{rf_heatmap}\end{center}
\vspace{-2mm}\begin{itemize}\small
\item Mỗi ô cho biết tỷ lệ hộ có giỏ trong ngày 538--544 và số hộ của nhóm tại mốc 537.
\item Nhóm dưới 30 hộ được tô xám vì tỷ lệ dễ thiếu ổn định; ranh giới các nhóm chỉ phục vụ mô tả.
\end{itemize}
\src{Notebook EDA, mục 11.2; bảng chi tiết vẫn giữ các ô nhỏ.}
\end{frame}

% EDA 13
\begin{frame}{Chi nhiều tiền có thể đi cùng việc mua nhiều lần}
\begin{center}\edashort{monetary_groups}\end{center}
\vspace{-2mm}\begin{itemize}\small
\item Tách hộ không có giỏ, hộ có giỏ nhưng doanh số không dương và bốn nhóm doanh số dương từ thấp đến cao.
\item Nhóm chi nhiều có thể cũng mua thường xuyên hơn; doanh số cao chưa có nghĩa lợi nhuận cao.
\end{itemize}
\src{Notebook phân tích dữ liệu, mục 11.3; bốn nhóm doanh số được chia tại mốc 537 để mô tả.}
\end{frame}

% EDA 14
\begin{frame}{Từ EDA đến quyết định sử dụng dữ liệu}
\small\renewcommand{\arraystretch}{1.18}
\begin{tabularx}{\textwidth}{>{\raggedright\arraybackslash}p{35mm}>{\raggedright\arraybackslash}X>{\raggedright\arraybackslash}X}
\toprule Điều quan sát được & Có thể biểu diễn tại mốc chốt & Cần tránh \\
\midrule
Lần mua gần nhất và nhịp mua & Số ngày từ lần mua gần nhất, số giỏ và số ngày mua trước mốc & Không dùng dữ liệu của 7 ngày tương lai. \\
Quy mô mua và giá trị & Số dòng sản phẩm, doanh số \texttt{SALES\_VALUE} & Không cộng \texttt{QUANTITY} giữa các nhóm hàng. \\
Hộ xuất hiện ở các thời điểm khác nhau & Thời gian kể từ lần đầu thấy hộ trong dữ liệu & Không coi đó là tuổi thực của khách hàng. \\
Giảm giá, nhiên liệu, thông tin hộ & So các hộ có nhịp mua tương tự; xem độ phủ và thời điểm có dữ liệu & Không suy ra quan hệ nhân quả hay khả năng dự báo chỉ từ EDA. \\
\bottomrule
\end{tabularx}
\vspace{3mm}
\src{Notebook phân tích dữ liệu, mục 12; các biến này chỉ là ứng viên, cần kiểm tra ở bước xây dựng mô hình.}
\end{frame}

% Kết thúc phần Data Understanding
\begin{frame}{Kết luận: đã hiểu dữ liệu, còn gì cần thận trọng?}
\begin{columns}[T,onlytextwidth]
\column{.49\textwidth}
\textbf{Đã xác lập}\vspace{2mm}
\begin{itemize}
\item Đã xác định mỗi dòng của tám bảng thể hiện gì và nối chúng ra sao.
\item Mua hàng nghĩa là có bất kỳ giỏ nào trong 7 ngày tiếp theo.
\item Lần mua gần nhất và nhịp mua có liên hệ rõ với việc mua tiếp.
\item Khi so nhóm nhiên liệu, giảm giá hay thông tin hộ, cần xét cả mức hoạt động.
\end{itemize}
\column{.48\textwidth}
\textbf{Giới hạn cần mang sang bước tiếp theo}\vspace{2mm}
\begin{itemize}
\item Chưa rõ đơn vị của \texttt{QUANTITY} ở mọi nhóm hàng.
\item Chưa biết ngày lịch và nguyên nhân ít giao dịch ở ngày 278, 643.
\item Chưa biết từ khi nào dữ liệu chiến dịch và thông tin hộ có sẵn.
\item Một hộ xuất hiện nhiều mốc; tỷ lệ theo lượt quan sát không phải tỷ lệ theo số hộ.
\end{itemize}
\end{columns}
\vspace{2mm}\src{Notebook \texttt{phan\_tich\_du\_lieu\_propensity.ipynb}; EDA hoàn thành trong phạm vi dữ liệu hiện có.}
\end{frame}

% Phụ lục: tên field và ví dụ được kiểm tra trên 8 CSV local.
\appendix

% 21
\begin{frame}{Phụ lục A1: Giao dịch — khóa và thời gian}
\small\renewcommand{\arraystretch}{1.18}
\begin{tabularx}{\textwidth}{>{\ttfamily}p{3.3cm}X}
\toprule Cột & Cách hiểu \\
\midrule
household\_key & Mã hộ gia đình có giao dịch. \\
BASKET\_ID & Mã giỏ hàng; nhiều dòng sản phẩm có thể cùng mã này. \\
DAY & Số thứ tự ngày trong dữ liệu; chưa xác định được ngày lịch tương ứng. \\
PRODUCT\_ID & Mã sản phẩm trên dòng giao dịch; dùng nối với bảng sản phẩm. \\
STORE\_ID & Mã cửa hàng ghi nhận giao dịch. \\
TRANS\_TIME & Giờ giao dịch dạng số HHMM; ví dụ 1631 là 16:31. \\
WEEK\_NO & Chỉ số tuần của giao dịch; không tự suy diễn là tuần lịch. \\
\bottomrule
\end{tabularx}
\vspace{2mm}\src{Mỗi dòng là một sản phẩm trong giỏ hàng; bảng giao dịch có 2.595.732 dòng.}
\end{frame}

% 22
\begin{frame}{Phụ lục A2: Giao dịch — số lượng và giá trị}
\small\renewcommand{\arraystretch}{1.12}
\begin{tabularx}{\textwidth}{>{\ttfamily}p{4.0cm}X}
\toprule Cột & Cách hiểu và lưu ý \\
\midrule
QUANTITY & Lượng sản phẩm mua; tài liệu chưa nêu rõ đơn vị riêng của nhiên liệu. \\
SALES\_VALUE & Số tiền nhà bán lẻ nhận được từ dòng sản phẩm; có thể khác số tiền khách trả. \\
RETAIL\_DISC & Giá trị điều chỉnh/chiết khấu bán lẻ; thường mang dấu âm trong dữ liệu. \\
COUPON\_DISC & Giá trị điều chỉnh/chiết khấu coupon được ghi nhận trên dòng. \\
COUPON\_MATCH\_DISC & Giá trị điều chỉnh coupon matching được ghi nhận trên dòng. \\
\bottomrule
\end{tabularx}
\vspace{1mm}\begin{block}{Cơ sở định nghĩa}
Tài liệu nguồn phân biệt số tiền nhà bán lẻ nhận được và số tiền khách trả khi dùng phiếu giảm giá của nhà sản xuất.
\end{block}
\src{The Complete Journey User Guide, trang 3.}
\end{frame}

% 23
\begin{frame}{Phụ lục B: Từ điển trường dữ liệu sản phẩm}
\small\renewcommand{\arraystretch}{1.16}
\begin{tabularx}{\textwidth}{>{\ttfamily}p{4.4cm}X}
\toprule Cột & Cách hiểu \\
\midrule
PRODUCT\_ID & Mã sản phẩm; dùng nối với giao dịch và phiếu giảm giá. \\
MANUFACTURER & Mã nhà sản xuất trong danh mục sản phẩm. \\
DEPARTMENT & Nhóm hàng lớn, ví dụ GROCERY và PRODUCE. \\
BRAND & Loại nhãn hàng: của nhà sản xuất (National) hoặc nhãn riêng (Private). \\
COMMODITY\_DESC & Mô tả nhóm hàng ở cấp commodity. \\
SUB\_COMMODITY\_DESC & Mô tả nhóm hàng chi tiết hơn dưới commodity. \\
CURR\_SIZE\_OF\_PRODUCT & Mô tả quy cách/kích cỡ sản phẩm; ví dụ \texttt{22 LB}. \\
\bottomrule
\end{tabularx}
\vspace{1mm}\src{92.353 dòng; PRODUCT\_ID không trùng trong catalog local.}
\end{frame}

% 24
\begin{frame}{Phụ lục C: Từ điển trường dữ liệu nhân khẩu học}
\scriptsize\renewcommand{\arraystretch}{1.15}
\begin{tabularx}{\textwidth}{>{\ttfamily}p{4.2cm}X}
\toprule Cột & Giá trị và cách hiểu từ dữ liệu hiện có \\
\midrule
household\_key & Mã hộ gia đình; dùng nối với giao dịch. \\
classification\_1 & Nhóm được mã hóa bằng nhãn \texttt{Age Group1}--\texttt{Age Group6}. \\
classification\_2 & Nhóm phân loại mã hóa; các giá trị gồm \texttt{X}, \texttt{Y}, \texttt{Z}. \\
classification\_3 & Nhóm phân loại mã hóa bằng các nhãn \texttt{Level}. \\
classification\_4 & Nhóm phân loại mã hóa dạng số; ý nghĩa cần xác nhận. \\
classification\_5 & Nhóm phân loại mã hóa bằng các nhãn \texttt{Group}. \\
HOMEOWNER\_DESC & Nhãn mô tả tình trạng nhà ở/sở hữu nhà, có giá trị \texttt{Unknown}. \\
KID\_CATEGORY\_DESC & Nhãn nhóm trẻ em trong hộ, gồm \texttt{None/Unknown}. \\
\bottomrule
\end{tabularx}
\vspace{1mm}\src{Chỉ 801 hộ có hồ sơ; chưa gán nghĩa cụ thể cho các mã classification khi thiếu tài liệu giải thích.}
\end{frame}

% 25
\begin{frame}{Phụ lục D: Từ điển trường phân bổ chiến dịch}
\small\renewcommand{\arraystretch}{1.3}
\begin{tabularx}{\textwidth}{>{\ttfamily}p{3.5cm}X}
\toprule Cột & Cách hiểu \\
\midrule
DESCRIPTION & Loại chiến dịch ghi trong bảng, ví dụ \texttt{TypeA}. \\
household\_key & Mã hộ được gắn với chiến dịch. \\
CAMPAIGN & Mã chiến dịch; dùng nối với \texttt{campaign\_desc}. \\
\bottomrule
\end{tabularx}
\vspace{5mm}\begin{block}{Đơn vị quan sát}
Mỗi dòng gắn một hộ với một chiến dịch. Điều đó chưa chứng minh hộ đã nhận, mở hoặc xem thông điệp.
\end{block}
\end{frame}

% 26
\begin{frame}{Phụ lục E: Từ điển trường mô tả chiến dịch}
\small\renewcommand{\arraystretch}{1.25}
\begin{tabularx}{\textwidth}{>{\ttfamily}p{3.5cm}X}
\toprule Cột & Cách hiểu \\
\midrule
DESCRIPTION & Loại chiến dịch, ví dụ \texttt{TypeB}. \\
CAMPAIGN & Mã chiến dịch; dùng nối với hộ được gắn chiến dịch hoặc lần dùng phiếu. \\
START\_DAY & Số thứ tự ngày bắt đầu chiến dịch. \\
END\_DAY & Số thứ tự ngày kết thúc chiến dịch. \\
\bottomrule
\end{tabularx}
\vspace{5mm}\begin{block}{Ý nghĩa phân tích}
Hai cột ngày giúp xác định thời gian chiến dịch và kiểm tra ngày dùng phiếu có nằm trong khoảng này không.
\end{block}
\end{frame}

% 27
\begin{frame}{Phụ lục F: Từ điển trường dữ liệu coupon}
\small\renewcommand{\arraystretch}{1.3}
\begin{tabularx}{\textwidth}{>{\ttfamily}p{3.5cm}X}
\toprule Cột & Cách hiểu \\
\midrule
COUPON\_UPC & Mã phiếu giảm giá; dùng nối với bảng ghi lần sử dụng. \\
PRODUCT\_ID & Mã sản phẩm được gắn với phiếu. \\
CAMPAIGN & Mã chiến dịch gắn với phiếu và sản phẩm. \\
\bottomrule
\end{tabularx}
\vspace{5mm}\begin{block}{Đơn vị quan sát và kết quả kiểm tra}
Mỗi dòng liên kết phiếu, sản phẩm và chiến dịch. Có 5.164 dòng giống hệt nhau trên cả ba cột.
\end{block}
\end{frame}

% 28
\begin{frame}{Phụ lục G: Từ điển trường sử dụng coupon}
\small\renewcommand{\arraystretch}{1.25}
\begin{tabularx}{\textwidth}{>{\ttfamily}p{3.5cm}X}
\toprule Cột & Cách hiểu \\
\midrule
household\_key & Mã hộ đã sử dụng phiếu giảm giá. \\
DAY & Số thứ tự ngày ghi nhận lần sử dụng phiếu. \\
COUPON\_UPC & Mã phiếu đã dùng; dùng nối với bảng phiếu giảm giá. \\
CAMPAIGN & Mã chiến dịch liên quan đến lần dùng phiếu. \\
\bottomrule
\end{tabularx}
\vspace{5mm}\begin{block}{Đơn vị quan sát}
Mỗi dòng ghi một lần hộ sử dụng phiếu trong một chiến dịch vào một ngày được ghi nhận.
\end{block}
\end{frame}

% 29
\begin{frame}{Phụ lục H: Từ điển trường trưng bày và mailer}
\small\renewcommand{\arraystretch}{1.23}
\begin{tabularx}{\textwidth}{>{\ttfamily}p{3.5cm}X}
\toprule Cột & Cách hiểu \\
\midrule
PRODUCT\_ID & Mã sản phẩm trong bản ghi trưng bày hoặc gửi quảng cáo. \\
STORE\_ID & Mã cửa hàng của bản ghi. \\
WEEK\_NO & Chỉ số tuần của bản ghi. \\
display & Mã hình thức trưng bày tại cửa hàng; chưa có bảng giải mã. \\
mailer & Mã hình thức gửi quảng cáo; chưa có bảng giải mã. \\
\bottomrule
\end{tabularx}
\vspace{3mm}\begin{block}{Đơn vị quan sát}
Mỗi dòng gắn một sản phẩm với một cửa hàng và một tuần, kèm mã trưng bày và gửi quảng cáo.
\end{block}
\src{36.786.524 dòng local; diễn giải mã display/mailer cần tài liệu gốc.}
\end{frame}
\end{document}
